\section*{Abstract}

Graph structure learning methods such as MTGNN and GTS infer relationships between time series jointly with a forecasting model, and report clear gains on traffic and weather benchmarks where signals are strong and near-stationary. We ask whether this transfers to financial returns, and whether gradient-based graph learning outperforms a derivative-free alternative. Using 20 semiconductor and agriculture companies over 1{,}806 trading days, we formulate next-day return prediction as node-level regression and compare MTGNN against VAR, an LSTM, CMA-ES searching the identical graph parameterization without gradients, and ablations replacing the learned graph with a fixed correlation matrix and with the identity.

The result is negative on every comparison. MTGNN attains test MAE $0.020948 \pm 0.000018$ against $0.020869$ for a zero forecast, with an information coefficient of $-0.0081$ (95\% CI $-0.050$ to $+0.034$) and a Diebold-Mariano $p$-value of $0.70$. CMA-ES is statistically indistinguishable from gradient learning ($p = 0.85$) while costing substantially more. Learned, correlation-based, CMA-ES-searched and identity adjacencies all fall within $2.0 \times 10^{-5}$ MAE, with the identity, which disables message passing entirely, nominally best. Of fourteen significance tests, only two reach the $5\%$ level and both identify models significantly \emph{worse} than predicting zero.

We explain the outcome through two mechanisms. Training on squared error against a near-unpredictable target drives every model to shrink its predictions to $5$--$18\%$ of the target's dispersion, so MAE rewards hedging rather than skill; the ordering of the error table is close to the reverse of the ordering of forecast quality. And because there is no cross-sectional signal to propagate, the loss is insensitive to the adjacency matrix, leaving the learned graph close to uninterpretable ($51.2\%$ of edge weight within-industry against a $47.4\%$ random baseline) while a fixed correlation graph recovers coherent economic clusters ($89.8\%$) and forecasts no better.

Our broader finding is methodological. We document three distinct ways this study could have reported a false positive, each caught only by the presence of a reference containing no information by construction. In low signal-to-noise domains, uninformative baselines are not a formality preceding the experiment; they are the instrument that makes it readable.
